\documentclass[11pt]{article}

% ---------- Packages ----------
\usepackage[a4paper,margin=1in]{geometry}
\usepackage{amsmath,amssymb,amsthm,mathtools}
\usepackage{hyperref}
\usepackage{tikz}
\usepackage{tikz-cd}
\usetikzlibrary{arrows.meta,positioning,cd}

% ---------- Theorem environments ----------
\newtheorem{theorem}{Theorem}
\newtheorem{definition}{Definition}
\newtheorem{lemma}{Lemma}
\newtheorem{proposition}{Proposition}
\newtheorem{remark}{Remark}
\newtheorem{example}{Example}

% ---------- Commands ----------
\newcommand{\Meas}{\mathbf{Meas}}
\newcommand{\Set}{\mathbf{Set}}
\newcommand{\Kl}{\mathrm{Kl}}
\newcommand{\Pcal}{\mathcal{P}}
\newcommand{\Xcal}{\mathcal{X}}
\newcommand{\Ycal}{\mathcal{Y}}
\newcommand{\Zcal}{\mathcal{Z}}
\newcommand{\Fcal}{\mathcal{F}}
\newcommand{\Bcal}{\mathcal{B}}
\newcommand{\Acal}{\mathcal{A}}
\newcommand{\Qcal}{\mathcal{Q}}
\newcommand{\id}{\mathrm{id}}
\newcommand{\op}{\mathrm{op}}

\title{
The Classical Construction of Probability:\\
A Categorical Reading
}

\author{Patrick S. Mahon}

\date{Working note --- \today}

\begin{document}

\maketitle

\begin{abstract}
This note develops the classical construction of probability theory in
categorical language, as preparation for identifying precisely where the
construction assumes a measure and what structure that assumption provides.
The goal is route-finding for the Observable Dynamics Program, which inverts
this construction.  We work through the category $\Meas$ of measurable spaces,
the Giry monad $\Pcal$ on $\Meas$, its Kleisli category $\Kl(\Pcal)$ of
Markov kernels, and the Kolmogorov extension theorem as a limit statement.
The assumed measure appears as a choice of monad algebra --- a free parameter
of the classical construction that the program seeks to derive.
\end{abstract}

\tableofcontents

%%%%%%%%%%%%%%%%%%%%%%%%%%%%%%%%%%%%%%%%%%%%%%%%%%%%%%%%%%%%%%%%%%%%%
\section{The category of measurable spaces}
%%%%%%%%%%%%%%%%%%%%%%%%%%%%%%%%%%%%%%%%%%%%%%%%%%%%%%%%%%%%%%%%%%%%%

\begin{definition}[The category $\Meas$]\label{def:meas}
The category $\Meas$ has:
\begin{itemize}
  \item \emph{Objects}: measurable spaces $(X, \Xcal)$ --- sets $X$ equipped
    with a $\sigma$-algebra $\Xcal \subseteq 2^X$.
  \item \emph{Morphisms}: measurable functions $f : (X, \Xcal) \to (Y, \Ycal)$
    --- functions $f : X \to Y$ such that $f^{-1}(B) \in \Xcal$ for every
    $B \in \Ycal$.
  \item \emph{Composition}: ordinary function composition.
  \item \emph{Identity}: the identity function $\id_X$.
\end{itemize}
\end{definition}

\begin{remark}[Products and limits in $\Meas$]
$\Meas$ has all small limits.  The product $(X \times Y,\, \Xcal \otimes \Ycal)$
carries the product $\sigma$-algebra generated by rectangles $A \times B$ with
$A \in \Xcal$, $B \in \Ycal$.  Countably infinite products
$\prod_{n} (X_n, \Xcal_n)$ carry the product $\sigma$-algebra generated by
cylinder sets --- preimages of measurable sets under finite coordinate
projections.  This is the measurable structure on the path space
$X^{\mathbb{Z}}$ used in the Observable Dynamics Program.
\end{remark}

\begin{remark}[What $\Meas$ does not capture]
The morphisms of $\Meas$ are deterministic: a measurable function sends each
point to exactly one point.  Stochastic processes, Markov chains, and
conditional distributions require a richer notion of morphism --- one that
sends points to \emph{distributions} rather than points.  This is provided
by the Kleisli category of the Giry monad, defined in Section~\ref{sec:kleisli}.
\end{remark}

%%%%%%%%%%%%%%%%%%%%%%%%%%%%%%%%%%%%%%%%%%%%%%%%%%%%%%%%%%%%%%%%%%%%%
\section{The Giry monad}\label{sec:giry}
%%%%%%%%%%%%%%%%%%%%%%%%%%%%%%%%%%%%%%%%%%%%%%%%%%%%%%%%%%%%%%%%%%%%%

\begin{definition}[The Giry functor]\label{def:giry-functor}
The \emph{Giry functor} $\Pcal : \Meas \to \Meas$ is defined as follows.
\begin{itemize}
  \item \emph{On objects}: $\Pcal(X, \Xcal)$ is the set of all probability
    measures on $(X, \Xcal)$, equipped with the \emph{evaluation
    $\sigma$-algebra} generated by the maps
    \[
      e_A : \Pcal(X) \to [0,1], \qquad \mu \mapsto \mu(A),
    \]
    for each $A \in \Xcal$.  We write $\Pcal X$ for $\Pcal(X, \Xcal)$ when
    the $\sigma$-algebra is clear.
  \item \emph{On morphisms}: given a measurable function $f : X \to Y$, the
    map $\Pcal f : \Pcal X \to \Pcal Y$ sends each measure $\mu$ to its
    \emph{pushforward} $f_\sharp \mu$:
    \[
      (\Pcal f)(\mu)(B) = (f_\sharp \mu)(B) = \mu(f^{-1}(B)), \qquad B \in \Ycal.
    \]
\end{itemize}
\end{definition}

\begin{proposition}[Functoriality of $\Pcal$]
$\Pcal$ is a functor: $\Pcal(\id_X) = \id_{\Pcal X}$ and
$\Pcal(g \circ f) = \Pcal g \circ \Pcal f$.
\end{proposition}

\begin{proof}
Both follow immediately from the definition of pushforward.
\end{proof}

\begin{definition}[The Giry monad]\label{def:giry-monad}
The \emph{Giry monad} is the triple $(\Pcal, \eta, \mu)$ where $\Pcal$ is the
Giry functor and:
\begin{itemize}
  \item \emph{Unit} $\eta : \id_{\Meas} \Rightarrow \Pcal$: the natural
    transformation whose component at $X$ is
    \[
      \eta_X : X \to \Pcal X, \qquad x \mapsto \delta_x,
    \]
    sending each point to the Dirac measure concentrated at that point.
  \item \emph{Multiplication} $\mu : \Pcal\Pcal \Rightarrow \Pcal$: the
    natural transformation whose component at $X$ is
    \[
      \mu_X : \Pcal(\Pcal X) \to \Pcal X,
      \qquad
      \mu_X(\Lambda)(A) = \int_{\Pcal X} \nu(A)\, \Lambda(d\nu),
    \]
    integrating a measure on measures down to a single measure.
\end{itemize}
These satisfy the monad axioms:
\[
  \mu_X \circ \eta_{\Pcal X} = \id_{\Pcal X}, \qquad
  \mu_X \circ \Pcal\eta_X = \id_{\Pcal X}, \qquad
  \mu_X \circ \mu_{\Pcal X} = \mu_X \circ \Pcal\mu_X.
\]
\end{definition}

\begin{remark}[What the multiplication means]
$\Pcal(\Pcal X)$ is the space of \emph{probability measures on probability
measures} --- a distribution over possible probability distributions.  This
arises naturally in Bayesian inference: a prior over models is a measure on
$\Pcal X$.  The multiplication $\mu_X$ collapses this to a single \emph{mixture
measure} by averaging.  Classically, a prior is a free choice of element of
$\Pcal(\Pcal X)$; the multiplication then yields the marginal distribution
on $X$.  The Observable Dynamics Program asks whether this choice can be
forced rather than freely made.
\end{remark}

\begin{remark}[Monad algebras]
A \emph{monad algebra} for $(\Pcal, \eta, \mu)$ is a measurable space $X$
together with a measurable map $\alpha : \Pcal X \to X$ satisfying:
\[
  \alpha \circ \eta_X = \id_X, \qquad \alpha \circ \mu_X = \alpha \circ \Pcal\alpha.
\]
Intuitively, $\alpha$ is an \emph{integration map} --- it evaluates a
distribution on $X$ to a point in $X$.  Monad algebras for the Giry monad
correspond to \emph{convex spaces}: spaces closed under convex combinations.
The classical probability space structure is a special case where $\alpha$
is evaluation of a fixed measure $P$.
\end{remark}

%%%%%%%%%%%%%%%%%%%%%%%%%%%%%%%%%%%%%%%%%%%%%%%%%%%%%%%%%%%%%%%%%%%%%
\section{The Kleisli category of Markov kernels}\label{sec:kleisli}
%%%%%%%%%%%%%%%%%%%%%%%%%%%%%%%%%%%%%%%%%%%%%%%%%%%%%%%%%%%%%%%%%%%%%

\begin{definition}[The Kleisli category $\Kl(\Pcal)$]\label{def:kleisli}
The \emph{Kleisli category} $\Kl(\Pcal)$ of the Giry monad has:
\begin{itemize}
  \item \emph{Objects}: measurable spaces $(X, \Xcal)$ --- the same objects
    as $\Meas$.
  \item \emph{Morphisms}: a morphism $K : X \to Y$ in $\Kl(\Pcal)$ is a
    \emph{Markov kernel} from $X$ to $Y$ --- a measurable function
    $K : X \to \Pcal Y$, i.e.\ a map that sends each point $x \in X$ to a
    probability measure $K(x, \cdot)$ on $Y$.
  \item \emph{Composition}: given $K : X \to Y$ and $L : Y \to Z$, their
    composite $L \circ K : X \to Z$ is the kernel
    \[
      (L \circ K)(x, C) = \int_Y L(y, C)\, K(x, dy), \qquad C \in \Zcal.
    \]
    This is the \emph{Chapman--Kolmogorov} composition.
  \item \emph{Identity}: $\id_X = \eta_X : X \to \Pcal X$, $x \mapsto \delta_x$.
\end{itemize}
\end{definition}

\begin{remark}[Deterministic maps embed into $\Kl(\Pcal)$]
Every measurable function $f : X \to Y$ in $\Meas$ gives a Markov kernel
$\hat{f} : X \to Y$ in $\Kl(\Pcal)$ via $\hat{f}(x, \cdot) = \delta_{f(x)}$.
This defines a functor $\Meas \to \Kl(\Pcal)$ that is the identity on objects
and sends deterministic maps to their Dirac-kernel counterparts.  In this sense
$\Kl(\Pcal)$ strictly extends $\Meas$ by allowing stochastic morphisms.
\end{remark}

\begin{remark}[Stochastic processes in $\Kl(\Pcal)$]
A (discrete-time) \emph{Markov chain} on $X$ is a morphism
$T : X \to X$ in $\Kl(\Pcal)$ --- a transition kernel.  A \emph{stochastic
process} is a diagram in $\Kl(\Pcal)$.  The Chapman--Kolmogorov identity
(that $n$-step transition kernels compose correctly) is exactly the
associativity of composition in $\Kl(\Pcal)$.  This makes $\Kl(\Pcal)$ the
natural ambient category for the dynamics studied in the Observable Dynamics
Program.
\end{remark}

\begin{example}[Refinement maps as Kleisli morphisms]
In the Observable Dynamics Program, a refinement map $\pi : O_j \to O_i$
between query outcome spaces is a deterministic measurable function.  It
embeds into $\Kl(\Pcal)$ as $\hat{\pi}(o, \cdot) = \delta_{\pi(o)}$.
The compatible marginals condition $(\pi_{ij})_\sharp \nu_j = \nu_i$ says
exactly that $\hat{\pi}_{ij}$ intertwines the marginals $\nu_j$ and $\nu_i$
--- a commutativity condition in $\Kl(\Pcal)$.
\end{example}

%%%%%%%%%%%%%%%%%%%%%%%%%%%%%%%%%%%%%%%%%%%%%%%%%%%%%%%%%%%%%%%%%%%%%
\section{Probability spaces as monad algebras}\label{sec:algebras}
%%%%%%%%%%%%%%%%%%%%%%%%%%%%%%%%%%%%%%%%%%%%%%%%%%%%%%%%%%%%%%%%%%%%%

\begin{definition}[Probability space as algebra]
A \emph{probability space} $(X, \Xcal, P)$ can be viewed as the measurable
space $(X, \Xcal)$ together with the choice of a \emph{point}
$P \in \Pcal X$ --- equivalently, a measurable map
\[
  P : 1 \to \Pcal X
\]
from the terminal object $1 = (\{\star\}, \{\emptyset, \{\star\}\})$ to
$\Pcal X$.  This is the \emph{free parameter} of the classical construction:
nothing in the structure of $(X, \Xcal)$ forces or determines $P$.
\end{definition}

\begin{remark}[The free choice]
This is the precise categorical location of the classical assumption.  The
measurable space $X$ carries no preferred measure.  Kolmogorov's axioms do
not say where the measure comes from --- they only say what properties it
must have once chosen.  The choice of $P$ is epistemologically a \emph{prior}
or \emph{model assumption}; it is not derived from the structure of
observations.

The Observable Dynamics Program asks: can $P$ be replaced by a \emph{canonical}
construction from observable data, so that the map $1 \to \Pcal X$ is not
freely chosen but forced by coherence constraints on a query system?
\end{remark}

\begin{remark}[Conditional distributions and disintegration]
Given a probability space $(X \times Y, P)$ and a measurable function
$f : X \times Y \to X$, a \emph{regular conditional distribution} is a
Markov kernel $K : X \to Y$ in $\Kl(\Pcal)$ such that
\[
  P(A \times B) = \int_A K(x, B)\, P_X(dx)
\]
where $P_X = (\Pcal f)(P)$ is the marginal on $X$.  This is the
\emph{disintegration} of $P$ over $f$.  Its existence (under standard Borel
hypotheses) is the Rokhlin disintegration theorem used in Paper~2 to
construct the predictive kernel.

Disintegration shows that the Kleisli morphism $K$ is \emph{derived from}
the algebra structure $P$ --- not freely chosen.  In the program, the
minimal predictive query reverses this: $K$ is derived from the query
structure, and $P$ is the algebra structure it forces.
\end{remark}

%%%%%%%%%%%%%%%%%%%%%%%%%%%%%%%%%%%%%%%%%%%%%%%%%%%%%%%%%%%%%%%%%%%%%
\section{Projective systems and the Kolmogorov extension theorem}
\label{sec:projective}
%%%%%%%%%%%%%%%%%%%%%%%%%%%%%%%%%%%%%%%%%%%%%%%%%%%%%%%%%%%%%%%%%%%%%

\begin{definition}[Projective system in $\Meas$]\label{def:proj-system}
A \emph{projective system} in $\Meas$ is a functor
$F : I^\op \to \Meas$ from a small directed category $I^\op$ to $\Meas$.
Concretely: a collection of measurable spaces $\{(X_i, \Xcal_i)\}_{i \in I}$
with measurable maps $\pi_{ij} : X_j \to X_i$ for each $i \leq j$, satisfying:
\[
  \pi_{ii} = \id_{X_i}, \qquad \pi_{ij} \circ \pi_{jk} = \pi_{ik}
  \quad (i \leq j \leq k).
\]
The \emph{projective limit} (inverse limit) is the measurable space
\[
  \varprojlim X_i = \left\{ (x_i)_{i \in I} \in \prod_i X_i \;\middle|\;
  \pi_{ij}(x_j) = x_i \text{ for all } i \leq j \right\}
\]
with the subspace $\sigma$-algebra from the product.
\end{definition}

\begin{remark}[Query systems as projective systems]
The query system $(\mathcal{Q}, \preceq)$ of the Observable Dynamics Program
is exactly a projective system in $\Meas$: the outcome spaces $\{O_i\}$ are
the objects and the refinement maps $\{\pi_{ij}\}$ are the morphisms.  The
realization space $\Omega = \varprojlim O_i$ is the projective limit.  The
Observable Dynamics Program therefore begins with the projective limit
already in place --- it is asking for the probability measure on this limit.
\end{remark}

\begin{definition}[Compatible family of measures]
Given a projective system $\{X_i, \pi_{ij}\}$, a \emph{compatible family}
of probability measures is a collection $\{\nu_i \in \Pcal X_i\}_{i \in I}$
satisfying
\[
  (\pi_{ij})_\sharp \nu_j = \nu_i \qquad \text{for all } i \leq j.
\]
This is a \emph{cocone} over the projective system in $\Kl(\Pcal)$: each
$\nu_i$ is consistent with the morphisms of the diagram.
\end{definition}

\begin{theorem}[Kolmogorov extension theorem --- classical form]
\label{thm:kolmogorov}
Let $\{(X_i, \Xcal_i), \pi_{ij}\}_{i \in I}$ be a projective system in
$\Meas$ where each $X_i$ is a Polish space with its Borel $\sigma$-algebra,
and let $\{\nu_i\}$ be a compatible family of probability measures.  Then
there exists a unique probability measure $P$ on the projective limit
$\Omega = \varprojlim X_i$ such that $(\mathrm{pr}_i)_\sharp P = \nu_i$ for
all $i \in I$, where $\mathrm{pr}_i : \Omega \to X_i$ is the projection.
\end{theorem}

\begin{remark}[Categorical reading]
The Kolmogorov extension theorem says that the functor
\[
  \Phi : \Pcal(\Omega) \to \left\{\text{compatible families } \{\nu_i\}\right\}
\]
sending $P$ to its marginals $\{(\mathrm{pr}_i)_\sharp P\}$ is an
\emph{isomorphism} (bijection) under Polish hypotheses.  The inverse of $\Phi$
is the extension construction.  Categorically, $\Phi$ is the counit of an
adjunction between the category of probability spaces over $\Omega$ and the
category of compatible marginal systems.

The Polish hypothesis is what forces the adjunction to exist.  Without it,
$\Phi$ is injective (uniqueness: the Observational Determination Theorem,
Paper~1) but not surjective (existence may fail).

The Observable Dynamics Program replaces the Polish hypothesis with
\emph{sequential upper-directedness} and \emph{realizability}, obtaining
existence by a different route --- a Carathéodory extension argument rather
than a topological compactness argument.
\end{remark}

%%%%%%%%%%%%%%%%%%%%%%%%%%%%%%%%%%%%%%%%%%%%%%%%%%%%%%%%%%%%%%%%%%%%%
\section{The assumed measure: a summary}
\label{sec:summary}
%%%%%%%%%%%%%%%%%%%%%%%%%%%%%%%%%%%%%%%%%%%%%%%%%%%%%%%%%%%%%%%%%%%%%

We can now state precisely where the classical construction assumes a measure
and what categorical structure that assumption provides.

\begin{center}
\begin{tabular}{lll}
\hline
Step & Construction & Categorical object \\
\hline
1 & Sample space $(X, \Xcal)$ & Object of $\Meas$ \\
2 & Giry monad $\Pcal$ & Monad on $\Meas$ \\
3 & \textbf{Choose measure} $P \in \Pcal X$ & \textbf{Free point} $1 \to \Pcal X$ \\
4 & Pushforward along $f$ & Functorial action $\Pcal f$ \\
5 & Markov kernels & Morphisms in $\Kl(\Pcal)$ \\
6 & Compatible marginals & Cocone in $\Kl(\Pcal)$ \\
7 & Kolmogorov extension & Limit/colimit in $\Meas$ (Polish case) \\
\hline
\end{tabular}
\end{center}

Step~3 is the free parameter.  Steps~1--2 and 4--5 are structural (forced by
the category).  Steps~6--7 are consequences of Step~3 plus the structural
machinery.

The Observable Dynamics Program begins at Step~6 and asks whether Step~3 can
be derived as a consequence --- whether the compatible cocone in $\Kl(\Pcal)$
forces a canonical point $1 \to \Pcal\Omega$ without a free choice.  The
extension theorems of Paper~1 give conditions under which it can.

\begin{remark}[The adjunction statement of Paper~1]
The main theorem of Paper~1 can be stated as: under sequential
upper-directedness and realizability, the functor $\Phi$ above is an
isomorphism.  Making this into a precise adjunction statement --- identifying
the two categories and the unit/counit --- is the next step of the categorical
development.  See the companion note on the program's picture.
\end{remark}

%%%%%%%%%%%%%%%%%%%%%%%%%%%%%%%%%%%%%%%%%%%%%%%%%%%%%%%%%%%%%%%%%%%%%
\section{Mathlib availability}\label{sec:mathlib}
%%%%%%%%%%%%%%%%%%%%%%%%%%%%%%%%%%%%%%%%%%%%%%%%%%%%%%%%%%%%%%%%%%%%%

See \texttt{mathlib\_gaps.md} (this directory) for the consolidated audit of
what Mathlib~4 has and lacks for this picture --- Giry monad, Kleisli, Markov
kernels present; projective limits in $\Meas$ and the Kolmogorov extension
theorem absent (formalization gaps, not mathematical ones). The classical
construction here builds on \texttt{GiryMonad} and \texttt{Kleisli} and connects
to the in-house \texttt{QuerySystem.lean} via the embedding $\Meas \hookrightarrow
\Kl(\Pcal)$.

\end{document}
